% gerado por ferramentas/md2tex_enunciado.py a partir de questoes/q096.md (não editar à mão)
Uma instalação de radar, usada para rastrear aviões, emite ondas eletromagnéticas com velocidade estimada de 300 mil quilômetros por segundo. Um avião que saiu da rota foi localizado em um raio de 75 quilômetros do aeroporto onde fica o radar. O operador usou o tempo dispendido para enviar e receber a resposta com a finalidade de estimar a localização do avião.\par
\figq[0.28\linewidth]{q096_fig1.png}
\fonte{Disponível em: www.pasternack.com. Acesso em: 1 abr. 2026 (adaptado).}
O tempo, em segundo, medido foi de\par
\begin{altlinha}
\alth{A}{$2{,}0 \times 10^{3}$\,s.}%
\alth{B}{$4{,}0 \times 10^{0}$\,s.}%
\alth{C}{$5{,}0 \times 10^{-1}$\,s.}%
\alth{D}{$1{,}5 \times 10^{-3}$\,s.}%
\alth{E}{$5{,}0 \times 10^{-4}$\,s.}
\end{altlinha}
